On prépare, à $\qty{25}{\degreeCelsius}$, une solution aqueuse
$(S_{A})$ d'acide butanoïque de concentration molaire
$C_{A}=\qty{2,0e-3}{\M}$ et de volume $V_{A}=\qty{1,0}{\liter}$. La
mesure du $\mathrm{pH}$ de la solution $(S_{A})$ a donné la valeur
$\mathrm{pH}=3,76$.
\begin{enumerate}
    \item Écrire l'équation chimique modélisant la réaction de l'acide
          butanoïque avec l'eau.
    \item Dresser le tableau d'avancement de la réaction en utilisant
          les grandeurs $C_{A}$, $V_{A}$, l'avancement $x$ et
          l'avancement $x_{\eq}$ à l'état d'équilibre du système
          chimique.
    \item Déterminer la valeur de l'avancement maximal $x_{\max}$.
    \item Vérifier que la valeur de l'avancement à l'état d'équilibre
          est $x_{\eq}=\qty{1,74e-4}{\mole}$.
    \item Calculer la valeur du taux d'avancement final $\tau$. Que
          peut-on déduire~?
    \item Calculer la valeur de la constante d'équilibre $\mathrm{K}$
          associée à l'équation de cette réaction.
    \item Recopier sur votre copie le numéro de la question et écrire
          la lettre correspondante à la proposition vraie.

          Dans les conditions de l'expérience, la constante
          d'équilibre $\mathrm{K}$ associée à l'équation de la
          réaction~:
          \begin{minipage}{\linewidth}
              \begin{table}[H]
                  \centering
                  \setlength{\arrayrulewidth}{0.8pt}
                  \renewcommand{\arraystretch}{1.3}
                  \begin{tabularx}{\textwidth}{|
                        >{\arraybackslash\centering}p{7mm}|L|}
                      \hline 
                      A &  dépend à la fois de la composition initiale du
                            système chimique et de la température. \\
                      \hline 
                      B &  dépend uniquement de la composition initiale du
                            système chimique. \\
                      \hline 
                      C &  dépend uniquement du $\mathrm{pH}$ de la solution. \\
                      \hline 
                      D &  dépend uniquement de la température du système
                            chimique. \\
                      \hline 
                  \end{tabularx}
              \end{table}
          \end{minipage}
          
    \item Calculer la valeur du
          $\mathrm{pK_{A}}(\ce{C3H7CO2H_{(aq)}/C3H7CO$_2^-$_{(aq)}})$.
\end{enumerate}

